% --- Archon physics formalization source begin ---
% archon:physics
% archon:covers PhyXMiniProblems/problem_phyx_mini_0971.lean
% archon:source-report reports/phyx_mini/problem_phyx_mini_0971.source.json

\chapter{Physics problem phyx\_mini\_0971}
\label{ch:PhyXMiniProblems_problem_phyx_mini_0971}

\paragraph{Problem source.}
\#\# Physical scenario

Two long, straight conducting wires with linear mass density \textbackslash{}( \textbackslash{}lambda \textbackslash{}) are suspended from cords such that they are each horizontal, parallel to each other, and a distance \textbackslash{}( d \textbackslash{}) apart. The back ends of the wires are connected to each other by a slack, low-resistance connecting wire. A charged capacitor (capacitance \textbackslash{}( C \textbackslash{})) is now added to the system; the positive plate of the capacitor (initial charge \textbackslash{}( +Q\_0 \textbackslash{})) is connected to the front end of one of the wires, and the negative plate of the capacitor (initial charge \textbackslash{}( -Q\_0 \textbackslash{})) is connected to the front end of the other wire. Both of these connections are also made by slack, low-resistance wires. When the connection is made, the wires are pushed aside by the repulsive force between the wires, and each wire has an initial horizontal velocity of magnitude \textbackslash{}( v\_0 \textbackslash{}). Assume that the time constant for the capacitor to discharge is negligible compared to the time it takes for any appreciable displacement in the position of the wires to occur.

\#\# Question

To what height \textbackslash{}( h \textbackslash{}) will each wire rise as a result of the circuit connection?

\#\# Answer choices

- A: A: \textbackslash{}frac\{1\}\{g\} \textbackslash{}left( \textbackslash{}frac\{\textbackslash{}mu\_0 Q\_0\textasciicircum{}2\}\{4 \textbackslash{}pi \textbackslash{}lambda RC d\} \textbackslash{}right)

- B: B: \textbackslash{}frac\{1\}\{2g\} \textbackslash{}left( \textbackslash{}frac\{\textbackslash{}mu\_0 Q\_0\textasciicircum{}2\}\{4 \textbackslash{}pi \textbackslash{}lambda RC d\} \textbackslash{}right)\textasciicircum{}2

- C: C: \textbackslash{}frac\{1\}\{2g\} \textbackslash{}left( \textbackslash{}frac\{\textbackslash{}mu\_0 Q\_0\textasciicircum{}2\}\{2 \textbackslash{}pi \textbackslash{}lambda RC d\} \textbackslash{}right)\textasciicircum{}2

- D: D: \textbackslash{}frac\{1\}\{2g\} \textbackslash{}left( \textbackslash{}frac\{\textbackslash{}mu\_0 Q\_0\textasciicircum{}2\}\{4 \textbackslash{}pi \textbackslash{}lambda RC d\textasciicircum{}2\} \textbackslash{}right)\textasciicircum{}2

\#\# Figure caption (auxiliary; use the image as primary evidence)

The image depicts an experimental setup involving two parallel rods or wires. 

- There are two rods positioned horizontally and parallel to each other. 
- The upper rod is depicted in yellow and connected to a circuit on the left side.
- The circuit includes a capacitor labeled "C" and is connected with positive and negative terminals.
- The lower rod is in gray.
- A vertical line indicates the distance between the two rods, labeled as "d".
- The rods are part of an electrical setup, possibly to demonstrate charging or interaction between electrical components. 

No additional background scene is present.

\#\# Dataset metadata

Electromagnetic Induction; Physical Model Grounding Reasoning, Multi-Formula Reasoning

\paragraph{Recorded answer/context.}
B: B: \textbackslash{}frac\{1\}\{2g\} \textbackslash{}left( \textbackslash{}frac\{\textbackslash{}mu\_0 Q\_0\textasciicircum{}2\}\{4 \textbackslash{}pi \textbackslash{}lambda RC d\} \textbackslash{}right)\textasciicircum{}2

\paragraph{Figure/image path.}
/root/proposal\_for\_physic/science-mango/phyx\_mini\_run/phyx\_data/test\_image/971.png

\paragraph{Formalization target.}
create a compiling Lean file with sorry bodies at `PhyXMiniProblems/problem\_phyx\_mini\_0971.lean`. The Lean declarations must preserve the physical quantities, dimensions or dimensional roles, figure labels, governing-law hypotheses, and final relation expressed by this problem.
Use Mathlib/Physlib names found through LeanExplore where available. If a physics API is missing, introduce faithful local abstractions rather than scalar placeholder aliases.

\begin{theorem}[Physics formalization target]
\label{thm:physics:phyx_mini_0971:target}
\lean{PhyXMiniProblems.ProblemPhyXMini0971.problem_phyx_mini_0971}
\uses{def:physics:phyx-mini-0971:phyxminiproblems-problemphyxmini0971-coherentsivalue, def:physics:phyx-mini-0971:phyxminiproblems-problemphyxmini0971-parallelwirecapacitorsetup, def:physics:phyx-mini-0971:phyxminiproblems-problemphyxmini0971-matcheswrittenparallelwirescenario, def:physics:phyx-mini-0971:phyxminiproblems-problemphyxmini0971-matchessuppliedparallelwirefigure, def:physics:phyx-mini-0971:phyxminiproblems-problemphyxmini0971-hasphysicalparallelwireparameters, def:physics:phyx-mini-0971:phyxminiproblems-problemphyxmini0971-satisfiesnegligibledisplacementasymptotics, def:physics:phyx-mini-0971:phyxminiproblems-problemphyxmini0971-satisfiesidealrcdischargelaw, def:physics:phyx-mini-0971:phyxminiproblems-problemphyxmini0971-satisfiesparallelwiremagneticforcelaw, def:physics:phyx-mini-0971:phyxminiproblems-problemphyxmini0971-satisfiesimpulsemomentumlaw, def:physics:phyx-mini-0971:phyxminiproblems-problemphyxmini0971-satisfiesballisticriseenergylaw, def:physics:phyx-mini-0971:phyxminiproblems-problemphyxmini0971-recordeddatasetanswer, def:physics:phyx-mini-0971:phyxminiproblems-problemphyxmini0971-answermatcheslimitingriseheight}
The assigned autoformalize agent should translate this physics problem into Lean declarations in the covered file, with theorem and lemma proof bodies written as `by sorry`.
\end{theorem}
\begin{proof}
This is an autoformalization task, not a proof task. Produce statements that can later be proved by the physics prover without weakening the physical model.
\end{proof}

% --- Archon declaration topology begin ---
\section{Declaration topology}
\begin{definition}[electric Current Dimension]
  \label{def:physics:phyx-mini-0971:phyxminiproblems-problemphyxmini0971-electriccurrentdimension}
  \lean{PhyXMiniProblems.ProblemPhyXMini0971.electricCurrentDimension}
  Electric current has dimension charge per time
\end{definition}
\begin{proof}
  This is the declared model definition of electric Current Dimension.
\end{proof}
\begin{definition}[electric Potential Dimension]
  \label{def:physics:phyx-mini-0971:phyxminiproblems-problemphyxmini0971-electricpotentialdimension}
  \lean{PhyXMiniProblems.ProblemPhyXMini0971.electricPotentialDimension}
  Electric potential has dimension energy per charge
\end{definition}
\begin{proof}
  This is the declared model definition of electric Potential Dimension.
\end{proof}
\begin{definition}[capacitance Dimension]
  \label{def:physics:phyx-mini-0971:phyxminiproblems-problemphyxmini0971-capacitancedimension}
  \lean{PhyXMiniProblems.ProblemPhyXMini0971.capacitanceDimension}
  \uses{def:physics:phyx-mini-0971:phyxminiproblems-problemphyxmini0971-electricpotentialdimension}
  Capacitance has dimension charge per electric potential
\end{definition}
\begin{proof}
  This is the declared model definition of capacitance Dimension.
\end{proof}
\begin{definition}[resistance Dimension]
  \label{def:physics:phyx-mini-0971:phyxminiproblems-problemphyxmini0971-resistancedimension}
  \lean{PhyXMiniProblems.ProblemPhyXMini0971.resistanceDimension}
  \uses{def:physics:phyx-mini-0971:phyxminiproblems-problemphyxmini0971-electricpotentialdimension}
  Resistance has dimension electric potential times time per charge
\end{definition}
\begin{proof}
  This is the declared model definition of resistance Dimension.
\end{proof}
\begin{definition}[linear Mass Density Dimension]
  \label{def:physics:phyx-mini-0971:phyxminiproblems-problemphyxmini0971-linearmassdensitydimension}
  \lean{PhyXMiniProblems.ProblemPhyXMini0971.linearMassDensityDimension}
  Linear mass density has dimension mass per length
\end{definition}
\begin{proof}
  This is the declared model definition of linear Mass Density Dimension.
\end{proof}
\begin{definition}[acceleration Dimension]
  \label{def:physics:phyx-mini-0971:phyxminiproblems-problemphyxmini0971-accelerationdimension}
  \lean{PhyXMiniProblems.ProblemPhyXMini0971.accelerationDimension}
  Acceleration has dimension length per time squared
\end{definition}
\begin{proof}
  This is the declared model definition of acceleration Dimension.
\end{proof}
\begin{definition}[magnetic Permeability Dimension]
  \label{def:physics:phyx-mini-0971:phyxminiproblems-problemphyxmini0971-magneticpermeabilitydimension}
  \lean{PhyXMiniProblems.ProblemPhyXMini0971.magneticPermeabilityDimension}
  Vacuum magnetic permeability has dimension mass length per charge squared
\end{definition}
\begin{proof}
  This is the declared model definition of magnetic Permeability Dimension.
\end{proof}
\begin{definition}[force Per Length Dimension]
  \label{def:physics:phyx-mini-0971:phyxminiproblems-problemphyxmini0971-forceperlengthdimension}
  \lean{PhyXMiniProblems.ProblemPhyXMini0971.forcePerLengthDimension}
  Force per unit length has dimension mass per time squared
\end{definition}
\begin{proof}
  This is the declared model definition of force Per Length Dimension.
\end{proof}
\begin{definition}[impulse Per Length Dimension]
  \label{def:physics:phyx-mini-0971:phyxminiproblems-problemphyxmini0971-impulseperlengthdimension}
  \lean{PhyXMiniProblems.ProblemPhyXMini0971.impulsePerLengthDimension}
  Impulse per unit length has dimension mass per time
\end{definition}
\begin{proof}
  This is the declared model definition of impulse Per Length Dimension.
\end{proof}
\begin{definition}[Magnitude]
  \label{def:physics:phyx-mini-0971:phyxminiproblems-problemphyxmini0971-magnitude}
  \lean{PhyXMiniProblems.ProblemPhyXMini0971.Magnitude}
  A nonnegative, unit-independent magnitude carrying dimension d
\end{definition}
\begin{proof}
  This is the declared model definition of Magnitude.
\end{proof}
\begin{definition}[Length Magnitude]
  \label{def:physics:phyx-mini-0971:phyxminiproblems-problemphyxmini0971-lengthmagnitude}
  \lean{PhyXMiniProblems.ProblemPhyXMini0971.LengthMagnitude}
  \uses{def:physics:phyx-mini-0971:phyxminiproblems-problemphyxmini0971-magnitude}
  This abbrev records the Length Magnitude component of the problem-specific physical model
\end{definition}
\begin{proof}
  This is the declared model definition of Length Magnitude.
\end{proof}
\begin{definition}[Time Magnitude]
  \label{def:physics:phyx-mini-0971:phyxminiproblems-problemphyxmini0971-timemagnitude}
  \lean{PhyXMiniProblems.ProblemPhyXMini0971.TimeMagnitude}
  \uses{def:physics:phyx-mini-0971:phyxminiproblems-problemphyxmini0971-magnitude}
  This abbrev records the Time Magnitude component of the problem-specific physical model
\end{definition}
\begin{proof}
  This is the declared model definition of Time Magnitude.
\end{proof}
\begin{definition}[Charge Magnitude]
  \label{def:physics:phyx-mini-0971:phyxminiproblems-problemphyxmini0971-chargemagnitude}
  \lean{PhyXMiniProblems.ProblemPhyXMini0971.ChargeMagnitude}
  \uses{def:physics:phyx-mini-0971:phyxminiproblems-problemphyxmini0971-magnitude}
  This abbrev records the Charge Magnitude component of the problem-specific physical model
\end{definition}
\begin{proof}
  This is the declared model definition of Charge Magnitude.
\end{proof}
\begin{definition}[Current Magnitude]
  \label{def:physics:phyx-mini-0971:phyxminiproblems-problemphyxmini0971-currentmagnitude}
  \lean{PhyXMiniProblems.ProblemPhyXMini0971.CurrentMagnitude}
  \uses{def:physics:phyx-mini-0971:phyxminiproblems-problemphyxmini0971-electriccurrentdimension, def:physics:phyx-mini-0971:phyxminiproblems-problemphyxmini0971-magnitude}
  This abbrev records the Current Magnitude component of the problem-specific physical model
\end{definition}
\begin{proof}
  This is the declared model definition of Current Magnitude.
\end{proof}
\begin{definition}[Capacitance Magnitude]
  \label{def:physics:phyx-mini-0971:phyxminiproblems-problemphyxmini0971-capacitancemagnitude}
  \lean{PhyXMiniProblems.ProblemPhyXMini0971.CapacitanceMagnitude}
  \uses{def:physics:phyx-mini-0971:phyxminiproblems-problemphyxmini0971-capacitancedimension, def:physics:phyx-mini-0971:phyxminiproblems-problemphyxmini0971-magnitude}
  This abbrev records the Capacitance Magnitude component of the problem-specific physical model
\end{definition}
\begin{proof}
  This is the declared model definition of Capacitance Magnitude.
\end{proof}
\begin{definition}[Resistance Magnitude]
  \label{def:physics:phyx-mini-0971:phyxminiproblems-problemphyxmini0971-resistancemagnitude}
  \lean{PhyXMiniProblems.ProblemPhyXMini0971.ResistanceMagnitude}
  \uses{def:physics:phyx-mini-0971:phyxminiproblems-problemphyxmini0971-resistancedimension, def:physics:phyx-mini-0971:phyxminiproblems-problemphyxmini0971-magnitude}
  This abbrev records the Resistance Magnitude component of the problem-specific physical model
\end{definition}
\begin{proof}
  This is the declared model definition of Resistance Magnitude.
\end{proof}
\begin{definition}[Linear Mass Density Magnitude]
  \label{def:physics:phyx-mini-0971:phyxminiproblems-problemphyxmini0971-linearmassdensitymagnitude}
  \lean{PhyXMiniProblems.ProblemPhyXMini0971.LinearMassDensityMagnitude}
  \uses{def:physics:phyx-mini-0971:phyxminiproblems-problemphyxmini0971-linearmassdensitydimension, def:physics:phyx-mini-0971:phyxminiproblems-problemphyxmini0971-magnitude}
  This abbrev records the Linear Mass Density Magnitude component of the problem-specific physical model
\end{definition}
\begin{proof}
  This is the declared model definition of Linear Mass Density Magnitude.
\end{proof}
\begin{definition}[Acceleration Magnitude]
  \label{def:physics:phyx-mini-0971:phyxminiproblems-problemphyxmini0971-accelerationmagnitude}
  \lean{PhyXMiniProblems.ProblemPhyXMini0971.AccelerationMagnitude}
  \uses{def:physics:phyx-mini-0971:phyxminiproblems-problemphyxmini0971-accelerationdimension, def:physics:phyx-mini-0971:phyxminiproblems-problemphyxmini0971-magnitude}
  This abbrev records the Acceleration Magnitude component of the problem-specific physical model
\end{definition}
\begin{proof}
  This is the declared model definition of Acceleration Magnitude.
\end{proof}
\begin{definition}[Magnetic Permeability Magnitude]
  \label{def:physics:phyx-mini-0971:phyxminiproblems-problemphyxmini0971-magneticpermeabilitymagnitude}
  \lean{PhyXMiniProblems.ProblemPhyXMini0971.MagneticPermeabilityMagnitude}
  \uses{def:physics:phyx-mini-0971:phyxminiproblems-problemphyxmini0971-magneticpermeabilitydimension, def:physics:phyx-mini-0971:phyxminiproblems-problemphyxmini0971-magnitude}
  This abbrev records the Magnetic Permeability Magnitude component of the problem-specific physical model
\end{definition}
\begin{proof}
  This is the declared model definition of Magnetic Permeability Magnitude.
\end{proof}
\begin{definition}[Force Per Length Magnitude]
  \label{def:physics:phyx-mini-0971:phyxminiproblems-problemphyxmini0971-forceperlengthmagnitude}
  \lean{PhyXMiniProblems.ProblemPhyXMini0971.ForcePerLengthMagnitude}
  \uses{def:physics:phyx-mini-0971:phyxminiproblems-problemphyxmini0971-forceperlengthdimension, def:physics:phyx-mini-0971:phyxminiproblems-problemphyxmini0971-magnitude}
  This abbrev records the Force Per Length Magnitude component of the problem-specific physical model
\end{definition}
\begin{proof}
  This is the declared model definition of Force Per Length Magnitude.
\end{proof}
\begin{definition}[Impulse Per Length Magnitude]
  \label{def:physics:phyx-mini-0971:phyxminiproblems-problemphyxmini0971-impulseperlengthmagnitude}
  \lean{PhyXMiniProblems.ProblemPhyXMini0971.ImpulsePerLengthMagnitude}
  \uses{def:physics:phyx-mini-0971:phyxminiproblems-problemphyxmini0971-impulseperlengthdimension, def:physics:phyx-mini-0971:phyxminiproblems-problemphyxmini0971-magnitude}
  This abbrev records the Impulse Per Length Magnitude component of the problem-specific physical model
\end{definition}
\begin{proof}
  This is the declared model definition of Impulse Per Length Magnitude.
\end{proof}
\begin{definition}[Speed Magnitude]
  \label{def:physics:phyx-mini-0971:phyxminiproblems-problemphyxmini0971-speedmagnitude}
  \lean{PhyXMiniProblems.ProblemPhyXMini0971.SpeedMagnitude}
  This abbrev records the Speed Magnitude component of the problem-specific physical model
\end{definition}
\begin{proof}
  This is the declared model definition of Speed Magnitude.
\end{proof}
\begin{definition}[coherent SIValue]
  \label{def:physics:phyx-mini-0971:phyxminiproblems-problemphyxmini0971-coherentsivalue}
  \lean{PhyXMiniProblems.ProblemPhyXMini0971.coherentSIValue}
  \uses{def:physics:phyx-mini-0971:phyxminiproblems-problemphyxmini0971-magnitude}
  The real-valued coherent-SI readout of a nonnegative physical magnitude
\end{definition}
\begin{proof}
  This is the declared model definition of coherent SIValue.
\end{proof}
\begin{definition}[Wire Label]
  \label{def:physics:phyx-mini-0971:phyxminiproblems-problemphyxmini0971-wirelabel}
  \lean{PhyXMiniProblems.ProblemPhyXMini0971.WireLabel}
  Conductors, circuit topology, and raster labels
\end{definition}
\begin{proof}
  This is the declared model definition of Wire Label.
\end{proof}
\begin{definition}[Wire End]
  \label{def:physics:phyx-mini-0971:phyxminiproblems-problemphyxmini0971-wireend}
  \lean{PhyXMiniProblems.ProblemPhyXMini0971.WireEnd}
  This inductive records the Wire End component of the problem-specific physical model
\end{definition}
\begin{proof}
  This is the declared model definition of Wire End.
\end{proof}
\begin{definition}[Capacitor Plate]
  \label{def:physics:phyx-mini-0971:phyxminiproblems-problemphyxmini0971-capacitorplate}
  \lean{PhyXMiniProblems.ProblemPhyXMini0971.CapacitorPlate}
  This inductive records the Capacitor Plate component of the problem-specific physical model
\end{definition}
\begin{proof}
  This is the declared model definition of Capacitor Plate.
\end{proof}
\begin{definition}[Along Wire Direction]
  \label{def:physics:phyx-mini-0971:phyxminiproblems-problemphyxmini0971-alongwiredirection}
  \lean{PhyXMiniProblems.ProblemPhyXMini0971.AlongWireDirection}
  This inductive records the Along Wire Direction component of the problem-specific physical model
\end{definition}
\begin{proof}
  This is the declared model definition of Along Wire Direction.
\end{proof}
\begin{definition}[Lateral Direction]
  \label{def:physics:phyx-mini-0971:phyxminiproblems-problemphyxmini0971-lateraldirection}
  \lean{PhyXMiniProblems.ProblemPhyXMini0971.LateralDirection}
  This inductive records the Lateral Direction component of the problem-specific physical model
\end{definition}
\begin{proof}
  This is the declared model definition of Lateral Direction.
\end{proof}
\begin{definition}[Wire Model]
  \label{def:physics:phyx-mini-0971:phyxminiproblems-problemphyxmini0971-wiremodel}
  \lean{PhyXMiniProblems.ProblemPhyXMini0971.WireModel}
  This inductive records the Wire Model component of the problem-specific physical model
\end{definition}
\begin{proof}
  This is the declared model definition of Wire Model.
\end{proof}
\begin{definition}[Wire Color]
  \label{def:physics:phyx-mini-0971:phyxminiproblems-problemphyxmini0971-wirecolor}
  \lean{PhyXMiniProblems.ProblemPhyXMini0971.WireColor}
  This inductive records the Wire Color component of the problem-specific physical model
\end{definition}
\begin{proof}
  This is the declared model definition of Wire Color.
\end{proof}
\begin{definition}[Parallel Wire Capacitor Figure]
  \label{def:physics:phyx-mini-0971:phyxminiproblems-problemphyxmini0971-parallelwirecapacitorfigure}
  \lean{PhyXMiniProblems.ProblemPhyXMini0971.ParallelWireCapacitorFigure}
  \uses{def:physics:phyx-mini-0971:phyxminiproblems-problemphyxmini0971-wirelabel, def:physics:phyx-mini-0971:phyxminiproblems-problemphyxmini0971-wireend, def:physics:phyx-mini-0971:phyxminiproblems-problemphyxmini0971-capacitorplate, def:physics:phyx-mini-0971:phyxminiproblems-problemphyxmini0971-wirecolor}
  Literal information read from the primary raster 971.png
\end{definition}
\begin{proof}
  This is the declared model definition of Parallel Wire Capacitor Figure.
\end{proof}
\begin{definition}[Parallel Wire Capacitor Setup]
  \label{def:physics:phyx-mini-0971:phyxminiproblems-problemphyxmini0971-parallelwirecapacitorsetup}
  \lean{PhyXMiniProblems.ProblemPhyXMini0971.ParallelWireCapacitorSetup}
  \uses{def:physics:phyx-mini-0971:phyxminiproblems-problemphyxmini0971-lengthmagnitude, def:physics:phyx-mini-0971:phyxminiproblems-problemphyxmini0971-timemagnitude, def:physics:phyx-mini-0971:phyxminiproblems-problemphyxmini0971-chargemagnitude, def:physics:phyx-mini-0971:phyxminiproblems-problemphyxmini0971-currentmagnitude, def:physics:phyx-mini-0971:phyxminiproblems-problemphyxmini0971-capacitancemagnitude, def:physics:phyx-mini-0971:phyxminiproblems-problemphyxmini0971-resistancemagnitude, def:physics:phyx-mini-0971:phyxminiproblems-problemphyxmini0971-linearmassdensitymagnitude, def:physics:phyx-mini-0971:phyxminiproblems-problemphyxmini0971-accelerationmagnitude, def:physics:phyx-mini-0971:phyxminiproblems-problemphyxmini0971-magneticpermeabilitymagnitude, def:physics:phyx-mini-0971:phyxminiproblems-problemphyxmini0971-forceperlengthmagnitude, def:physics:phyx-mini-0971:phyxminiproblems-problemphyxmini0971-impulseperlengthmagnitude, def:physics:phyx-mini-0971:phyxminiproblems-problemphyxmini0971-speedmagnitude, def:physics:phyx-mini-0971:phyxminiproblems-problemphyxmini0971-wirelabel, def:physics:phyx-mini-0971:phyxminiproblems-problemphyxmini0971-wireend, def:physics:phyx-mini-0971:phyxminiproblems-problemphyxmini0971-capacitorplate, def:physics:phyx-mini-0971:phyxminiproblems-problemphyxmini0971-alongwiredirection, def:physics:phyx-mini-0971:phyxminiproblems-problemphyxmini0971-lateraldirection, def:physics:phyx-mini-0971:phyxminiproblems-problemphyxmini0971-wiremodel, def:physics:phyx-mini-0971:phyxminiproblems-problemphyxmini0971-parallelwirecapacitorfigure}
  Physical data and coherent-SI-time histories for a family of otherwise identical experiments. The first real argument of the epsilon-dependent fields is the positive discharge-to-motion timescale ratio. Rise height is an independent observable and is not defined by an answer expression
\end{definition}
\begin{proof}
  This is the declared model definition of Parallel Wire Capacitor Setup.
\end{proof}
\begin{definition}[Matches Written Parallel Wire Scenario]
  \label{def:physics:phyx-mini-0971:phyxminiproblems-problemphyxmini0971-matcheswrittenparallelwirescenario}
  \lean{PhyXMiniProblems.ProblemPhyXMini0971.MatchesWrittenParallelWireScenario}
  \uses{def:physics:phyx-mini-0971:phyxminiproblems-problemphyxmini0971-parallelwirecapacitorsetup}
  The physical arrangement stated in the prose, independently of the raster
\end{definition}
\begin{proof}
  This is the declared model definition of Matches Written Parallel Wire Scenario.
\end{proof}
\begin{definition}[Matches Supplied Parallel Wire Figure]
  \label{def:physics:phyx-mini-0971:phyxminiproblems-problemphyxmini0971-matchessuppliedparallelwirefigure}
  \lean{PhyXMiniProblems.ProblemPhyXMini0971.MatchesSuppliedParallelWireFigure}
  \uses{def:physics:phyx-mini-0971:phyxminiproblems-problemphyxmini0971-parallelwirecapacitorsetup}
  Primary-image evidence. It contains no height or answer-choice claim
\end{definition}
\begin{proof}
  This is the declared model definition of Matches Supplied Parallel Wire Figure.
\end{proof}
\begin{definition}[Has Physical Parallel Wire Parameters]
  \label{def:physics:phyx-mini-0971:phyxminiproblems-problemphyxmini0971-hasphysicalparallelwireparameters}
  \lean{PhyXMiniProblems.ProblemPhyXMini0971.HasPhysicalParallelWireParameters}
  \uses{def:physics:phyx-mini-0971:phyxminiproblems-problemphyxmini0971-coherentsivalue, def:physics:phyx-mini-0971:phyxminiproblems-problemphyxmini0971-parallelwirecapacitorsetup}
  Positivity conditions for all dimensional parameters used below
\end{definition}
\begin{proof}
  This is the declared model definition of Has Physical Parallel Wire Parameters.
\end{proof}
\begin{definition}[fixed Separation Force Per Length SI]
  \label{def:physics:phyx-mini-0971:phyxminiproblems-problemphyxmini0971-fixedseparationforceperlengthsi}
  \lean{PhyXMiniProblems.ProblemPhyXMini0971.fixedSeparationForcePerLengthSI}
  \uses{def:physics:phyx-mini-0971:phyxminiproblems-problemphyxmini0971-coherentsivalue, def:physics:phyx-mini-0971:phyxminiproblems-problemphyxmini0971-parallelwirecapacitorsetup}
  The force per unit length predicted by the long-wire law when the separation is held at its initial value. This is a reference integrand for the asymptotic error contract, not the actual force and not a height formula
\end{definition}
\begin{proof}
  This is the declared model definition of fixed Separation Force Per Length SI.
\end{proof}
\begin{definition}[Satisfies Negligible Displacement Asymptotics]
  \label{def:physics:phyx-mini-0971:phyxminiproblems-problemphyxmini0971-satisfiesnegligibledisplacementasymptotics}
  \lean{PhyXMiniProblems.ProblemPhyXMini0971.SatisfiesNegligibleDisplacementAsymptotics}
  \uses{def:physics:phyx-mini-0971:phyxminiproblems-problemphyxmini0971-coherentsivalue, def:physics:phyx-mini-0971:phyxminiproblems-problemphyxmini0971-parallelwirecapacitorsetup, def:physics:phyx-mini-0971:phyxminiproblems-problemphyxmini0971-fixedseparationforceperlengthsi}
  A valid small-timescale contract replacing the former globalized fixed-separation equality. On every fixed multiple of the RC time, the separation approaches d. The L1 force error over the whole positive time ray is finite and tends to zero; this supplies the domination/tail control needed to pass from pointwise geometry to impulse. No speed or height conclusion is assumed here
\end{definition}
\begin{proof}
  This is the declared model definition of Satisfies Negligible Displacement Asymptotics.
\end{proof}
\begin{definition}[Satisfies Ideal RCDischarge Law]
  \label{def:physics:phyx-mini-0971:phyxminiproblems-problemphyxmini0971-satisfiesidealrcdischargelaw}
  \lean{PhyXMiniProblems.ProblemPhyXMini0971.SatisfiesIdealRCDischargeLaw}
  \uses{def:physics:phyx-mini-0971:phyxminiproblems-problemphyxmini0971-coherentsivalue, def:physics:phyx-mini-0971:phyxminiproblems-problemphyxmini0971-parallelwirecapacitorsetup}
  Exact ideal-RC discharge law, including tau = R C
\end{definition}
\begin{proof}
  This is the declared model definition of Satisfies Ideal RCDischarge Law.
\end{proof}
\begin{definition}[Satisfies Parallel Wire Magnetic Force Law]
  \label{def:physics:phyx-mini-0971:phyxminiproblems-problemphyxmini0971-satisfiesparallelwiremagneticforcelaw}
  \lean{PhyXMiniProblems.ProblemPhyXMini0971.SatisfiesParallelWireMagneticForceLaw}
  \uses{def:physics:phyx-mini-0971:phyxminiproblems-problemphyxmini0971-coherentsivalue, def:physics:phyx-mini-0971:phyxminiproblems-problemphyxmini0971-parallelwirecapacitorsetup}
  Exact long-parallel-wire magnetic force law at the actual separation
\end{definition}
\begin{proof}
  This is the declared model definition of Satisfies Parallel Wire Magnetic Force Law.
\end{proof}
\begin{definition}[Satisfies Impulse Momentum Law]
  \label{def:physics:phyx-mini-0971:phyxminiproblems-problemphyxmini0971-satisfiesimpulsemomentumlaw}
  \lean{PhyXMiniProblems.ProblemPhyXMini0971.SatisfiesImpulseMomentumLaw}
  \uses{def:physics:phyx-mini-0971:phyxminiproblems-problemphyxmini0971-coherentsivalue, def:physics:phyx-mini-0971:phyxminiproblems-problemphyxmini0971-parallelwirecapacitorsetup}
  Impulse per unit length is integrated force and equals lambda times speed
\end{definition}
\begin{proof}
  This is the declared model definition of Satisfies Impulse Momentum Law.
\end{proof}
\begin{definition}[Satisfies Ballistic Rise Energy Law]
  \label{def:physics:phyx-mini-0971:phyxminiproblems-problemphyxmini0971-satisfiesballisticriseenergylaw}
  \lean{PhyXMiniProblems.ProblemPhyXMini0971.SatisfiesBallisticRiseEnergyLaw}
  \uses{def:physics:phyx-mini-0971:phyxminiproblems-problemphyxmini0971-coherentsivalue, def:physics:phyx-mini-0971:phyxminiproblems-problemphyxmini0971-parallelwirecapacitorsetup}
  Subsequent kinetic energy per unit length becomes gravitational energy
\end{definition}
\begin{proof}
  This is the declared model definition of Satisfies Ballistic Rise Energy Law.
\end{proof}
\begin{theorem}[horizontal kick speed tends to capacitor impulse]
  \label{thm:physics:phyx-mini-0971:phyxminiproblems-problemphyxmini0971-horizontal-kick-speed-tends-to-capacitor-impulse}
  \lean{PhyXMiniProblems.ProblemPhyXMini0971.horizontal_kick_speed_tends_to_capacitor_impulse}
  \uses{def:physics:phyx-mini-0971:phyxminiproblems-problemphyxmini0971-coherentsivalue, def:physics:phyx-mini-0971:phyxminiproblems-problemphyxmini0971-parallelwirecapacitorsetup, def:physics:phyx-mini-0971:phyxminiproblems-problemphyxmini0971-hasphysicalparallelwireparameters, def:physics:phyx-mini-0971:phyxminiproblems-problemphyxmini0971-satisfiesnegligibledisplacementasymptotics, def:physics:phyx-mini-0971:phyxminiproblems-problemphyxmini0971-satisfiesidealrcdischargelaw, def:physics:phyx-mini-0971:phyxminiproblems-problemphyxmini0971-satisfiesparallelwiremagneticforcelaw, def:physics:phyx-mini-0971:phyxminiproblems-problemphyxmini0971-satisfiesimpulsemomentumlaw}
  The small-timescale limit of the electromagnetic kick speed. This is an intermediate consequence of the governing laws, not a premise of the height theorem
\end{theorem}
\begin{proof}
  The conclusion follows from the governing relations and figure readouts named in the statement of horizontal kick speed tends to capacitor impulse.
\end{proof}
\begin{definition}[Answer Choice]
  \label{def:physics:phyx-mini-0971:phyxminiproblems-problemphyxmini0971-answerchoice}
  \lean{PhyXMiniProblems.ProblemPhyXMini0971.AnswerChoice}
  Labels of the four alternatives printed by the source dataset
\end{definition}
\begin{proof}
  This is the declared model definition of Answer Choice.
\end{proof}
\begin{definition}[displayed Expression]
  \label{def:physics:phyx-mini-0971:phyxminiproblems-problemphyxmini0971-answerchoice-displayedexpression}
  \lean{PhyXMiniProblems.ProblemPhyXMini0971.AnswerChoice.displayedExpression}
  \uses{def:physics:phyx-mini-0971:phyxminiproblems-problemphyxmini0971-coherentsivalue, def:physics:phyx-mini-0971:phyxminiproblems-problemphyxmini0971-parallelwirecapacitorsetup, def:physics:phyx-mini-0971:phyxminiproblems-problemphyxmini0971-answerchoice}
  Coherent-SI scalar expression printed for each source answer choice
\end{definition}
\begin{proof}
  This is the declared model definition of displayed Expression.
\end{proof}
\begin{definition}[recorded Dataset Answer]
  \label{def:physics:phyx-mini-0971:phyxminiproblems-problemphyxmini0971-recordeddatasetanswer}
  \lean{PhyXMiniProblems.ProblemPhyXMini0971.recordedDatasetAnswer}
  \uses{def:physics:phyx-mini-0971:phyxminiproblems-problemphyxmini0971-answerchoice}
  The answer label recorded in the dataset
\end{definition}
\begin{proof}
  This is the declared model definition of recorded Dataset Answer.
\end{proof}
\begin{definition}[Answer Matches Limiting Rise Height]
  \label{def:physics:phyx-mini-0971:phyxminiproblems-problemphyxmini0971-answermatcheslimitingriseheight}
  \lean{PhyXMiniProblems.ProblemPhyXMini0971.AnswerMatchesLimitingRiseHeight}
  \uses{def:physics:phyx-mini-0971:phyxminiproblems-problemphyxmini0971-coherentsivalue, def:physics:phyx-mini-0971:phyxminiproblems-problemphyxmini0971-parallelwirecapacitorsetup, def:physics:phyx-mini-0971:phyxminiproblems-problemphyxmini0971-answerchoice}
  A source answer equals the limiting rise height of both conductors
\end{definition}
\begin{proof}
  This is the declared model definition of Answer Matches Limiting Rise Height.
\end{proof}
% --- Archon declaration topology end ---
% --- Archon physics formalization source end ---
